\documentclass[11pt]{article}
\usepackage[utf8]{inputenc}
\usepackage[spanish]{babel}
\usepackage{amsmath,amssymb,amsthm}
\usepackage{graphicx}
\usepackage[margin=2.5cm]{geometry}
\usepackage{enumitem}

\title{Semana 4 -- Ejercicios}
\author{Tutorías -- Introducción al Cálculo}
\date{24/08--28/08/2026}

\begin{document}
\maketitle

\section{Ejercicios de la semana}

\begin{enumerate}[label=\arabic*.]
  \item Escriba $[-3,2] \cup \,]0,5[$ como unión de intervalos disjuntos.

  \item Calcule $\,]-\infty,4] \cap [1,\infty[\,$. (Respuesta: $[1,4]$.)

  \item Resuelva $5x+3 \le 2x-9$. (Respuesta: $\,]-\infty,-4]$.)

  \item Resuelva $-4x+1<9$. (Respuesta: $\,]-2,\infty[$.)

  \item Resuelva $x^2-5x+6<0$, indicando discriminante, raíces y el conjunto solución.

  \item Resuelva $2x^2+x+1>0$. (Indicación: calcule $\Delta$ primero.)

  \item Resuelva $-x^2+4>0$.

  \item Resuelva, usando la misma idea del Problema 2.3 del libro,
  $(x^2-2x-3)(x^2+1) \le 0$.

  \item Resuelva $x^2+4x+4 \ge 0$ (caso $\Delta=0$: analice qué pasa en la raíz doble).

  \item Determine todos los $x\in\mathbb{R}$ tales que $(2x-1)(x+3)(x^2+2)<0$.

  \item \textbf{(Nivel interrogación)} Resuelva la inecuación racional
  $\dfrac{x^2-5x+6}{x^2-9} \le 0$, indicando explícitamente los valores de $x$ excluidos del
  dominio y justificando el resultado con una tabla de signos completa.

  \item \textbf{(Nivel interrogación)} Determine para qué valores de $m\in\mathbb{R}$ la ecuación
  $x^2 - (m+2)x + (2m+1) = 0$ tiene dos raíces reales \textbf{distintas}, justificando con el
  signo del discriminante.

  \item \textbf{(Nivel interrogación)} Resuelva la inecuación
  $\dfrac{(x-1)(x+2)}{x^2-4} \ge 0$. \textbf{No cancele} el factor común $(x+2)$ sin antes
  considerar su efecto en el dominio: indique explícitamente qué le ocurre a la expresión en
  $x=-2$.
\end{enumerate}

\end{document}
